%% mg_deep.tex — Chapters 22–26 of Meaning Games (v2)
%% No preamble; \input from Meaning_Games.tex
%% These chapters form Part V: Deep Applications and Grand Synthesis

%% ================================================================
%%  CHAPTER 22: THE POLITICAL ECONOMY OF MEANING
%% ================================================================

\chapter{The Political Economy of Meaning}
\label{ch:political-economy}

\epigraph{``The limits of my language mean the limits of my world.''\\ \medskip
``A serious and good philosophical work could be written consisting
entirely of jokes.''}{Ludwig Wittgenstein, \textit{Tractatus} 5.6 and \textit{Culture and Value}}

Politics is, at its deepest level, a contest over meaning.
When a government censors a newspaper, it is not merely suppressing
information---it is strategically removing ideas from the public
ideatic space so that certain compositions become impossible.
When a demagogue coins a slogan, they are injecting a high-emergence
idea designed to reshape the resonance landscape of an entire polity.
When citizens deliberate in a town hall, they are playing a
cooperative meaning game whose equilibrium---if one exists---is
what Habermas called the ``ideal speech situation.''

This chapter formalises these phenomena.  We assume throughout that
the reader is familiar with the eight axioms of IDT (especially the
composition operation~$\compose$, the resonance function~$\rs$,
the emergence term~$\emerge$, and the void~$\void$) and with
the basic theory of meaning games developed in Chapters~2--7.

%% ---------------------------------------------------------------
\section{The Marketplace of Ideas}
\label{sec:marketplace}

The metaphor of a ``marketplace of ideas'' dates to Justice Oliver
Wendell Holmes's dissent in \textit{Abrams v.\ United States} (1919),
but it has never been given a precise mathematical formulation.
We provide one now.

\begin{definition}[Ideatic marketplace]
\label{def:marketplace}
An \textbf{ideatic marketplace} is a tuple
$\mathcal{M} = (N, \IS, \rs, (\mathcal{A}_i)_{i \in N}, u)$ where:
\begin{enumerate}[label=(\roman*)]
  \item $N = \{1,\ldots,n\}$ is a set of \textbf{citizens} (players);
  \item $(\IS, \compose, \rs, \emerge)$ is an ideatic space
    satisfying the IDT axioms;
  \item $\mathcal{A}_i \subseteq \IS$ is the \textbf{ideatic endowment}
    of citizen~$i$---the set of ideas available to them;
  \item $u_i \colon \IS \to \R$ is a utility function representing
    citizen~$i$'s preferences over ideas, with the constraint that
    $u_i(a) \ge u_i(\void)$ for all $a \in \mathcal{A}_i$.
\end{enumerate}
A \textbf{discourse} is a sequence of compositions
$\sigma = (a_{i_1}, a_{i_2}, \ldots, a_{i_T})$ where
$a_{i_t} \in \mathcal{A}_{i_t}$ and the resulting public idea is
$\Pi_\sigma = a_{i_1} \compose a_{i_2} \compose \cdots \compose a_{i_T}$.
\end{definition}

The marketplace of ideas hypothesis asserts, informally, that free
and open discourse leads to truth.  The next definition and theorem
make this precise---and show when it fails.

\begin{definition}[Truth-tracking]
\label{def:truth-tracking}
Let $a^* \in \IS$ be a \textbf{ground truth} (an idea with maximal
resonance in some objective sense: $\rs(a^*,a^*) \ge \rs(a,a)$
for all $a \in \IS$).  A marketplace $\mathcal{M}$ is
\textbf{truth-tracking} if, for every Nash equilibrium discourse
$\sigma^*$, the public idea $\Pi_{\sigma^*}$ satisfies
\[
  \rs(\Pi_{\sigma^*},\, a^*) \;\ge\;
  (1-\delta)\,\rs(a^*,a^*),
\]
for some tolerance $\delta \in [0,1)$.
\end{definition}

\begin{theorem}[Marketplace failure under asymmetric emergence]
\label{thm:marketplace-failure}
Let $\mathcal{M}$ be an ideatic marketplace with $n \ge 2$ citizens.
Suppose there exist ideas $p \in \mathcal{A}_1$ (propaganda) and
$a^* \in \mathcal{A}_2$ (truth) such that:
\begin{enumerate}[label=(\alph*)]
  \item $\emerge(p, c, p \compose c) > \emerge(a^*, c, a^* \compose c)$
    for all $c \in \IS \setminus \{\void\}$;
  \item $u_1(p \compose c) > u_1(a^* \compose c)$ for all $c$.
\end{enumerate}
Then in every Nash equilibrium, citizen~$1$ plays $p$, and the
marketplace is not truth-tracking.
\end{theorem}

\begin{proof}
Condition~(a) ensures that propaganda~$p$ creates more emergence
than truth~$a^*$ when composed with any idea.  By the decomposition
axiom (Axiom~A5 of IDT), the weight of any discourse containing~$p$
exceeds the weight of the corresponding discourse containing~$a^*$:
\[
  w(p \compose c) = w(p) + \rs(c, p \compose c)
  + \emerge(p,c,p\compose c)
  > w(a^*) + \rs(c, a^* \compose c)
  + \emerge(a^*,c,a^*\compose c)
  = w(a^* \compose c).
\]
Condition~(b) ensures that citizen~$1$ strictly prefers any
discourse containing~$p$ to the corresponding discourse
containing~$a^*$.  Since this holds regardless of others'
strategies, playing~$p$ is a strictly dominant strategy for
citizen~$1$.  In every Nash equilibrium, citizen~$1$ plays~$p$.

To see that truth-tracking fails, note that if $\rs(p, a^*) < \rs(a^*,a^*)$,
then $\rs(\Pi_{\sigma^*}, a^*) \le \rs(p \compose c, a^*)
< \rs(a^*,a^*)$ for sufficiently small resonance between $p$ and $a^*$,
violating the truth-tracking condition for $\delta$ close to $0$.
\end{proof}

\begin{interpretation}
The Marketplace Failure Theorem formalises what media scholars
have long suspected: the ``marketplace of ideas'' can be systematically
corrupted when some participants possess ideas with higher emergence
than the truth.  Propaganda that ``goes viral'' (positive emergence
with every context) will dominate every equilibrium, regardless of
the truth's intrinsic resonance.  This is not a failure of
rationality---it is a structural consequence of the emergence
asymmetry.
\end{interpretation}

%% ---------------------------------------------------------------
\section{Propaganda Games}
\label{sec:propaganda}

We now develop a formal model of propaganda as a specific type
of meaning game with asymmetric information.

\begin{game}[The propaganda game]
\label{game:propaganda}
A \textbf{propaganda game} is a Bayesian meaning game
$\Gamma_P = (S, R, \IS, \rs, \Theta, p, u_S, u_R)$ where:
\begin{itemize}
  \item $S$ is the \textbf{propagandist} (sender),
    $R$ is the \textbf{public} (receiver);
  \item $\Theta = \{\theta_T, \theta_F\}$ is the state space
    (true state, false state);
  \item $p(\theta_T) = \pi$, $p(\theta_F) = 1-\pi$ is the common prior;
  \item $S$ observes $\theta$ and sends a message $m \in \IS$;
  \item $R$ observes $m$ and forms a posterior belief;
  \item $u_S(\theta, m) = \rs(m, a_F) + \beta\,\emerge(m, b_R, m \compose b_R)$,
    where $a_F$ is the propagandist's preferred narrative and $b_R$
    is the receiver's current belief;
  \item $u_R(\theta, m) = \rs(m, a_\theta)$---the receiver wants
    messages that resonate with the true state.
\end{itemize}
\end{game}

\begin{definition}[Propaganda effectiveness]
\label{def:propaganda-effectiveness}
The \textbf{propaganda effectiveness} of message $m$ against
belief $b$ is
\[
  \mathrm{PE}(m, b) \;=\;
  \frac{\rs(m \compose b, a_F)}{\rs(m \compose b, a_T)}
  \cdot
  \frac{\emerge(m, b, m \compose b)}{w(b)}.
\]
Propaganda is \textbf{successful} if $\mathrm{PE}(m,b) > 1$:
the composed belief resonates more with the false narrative
than with the truth, amplified by emergence.
\end{definition}

\begin{theorem}[Propaganda convergence]
\label{thm:propaganda-convergence}
In the repeated propaganda game with discount factor $\delta < 1$,
if the propagandist plays a stationary strategy $m^*$ with
$\mathrm{PE}(m^*, b_t) > 1$ for all $t$, then the receiver's
belief sequence $(b_t)_{t \ge 0}$ satisfies
\[
  \lim_{t \to \infty} \rs(b_t, a_F) = \rs(a_F, a_F).
\]
That is, the receiver's beliefs converge to the propagandist's
preferred narrative.
\end{theorem}

\begin{proof}
Define $d_t = \rs(b_t, a_F) / w(a_F)$ as the normalised alignment
with the false narrative.  After each round, the receiver's belief
updates to $b_{t+1} = b_t \compose m^*$.  By the decomposition axiom:
\begin{align*}
  \rs(b_{t+1}, a_F)
  &= \rs(b_t \compose m^*, a_F) \\
  &\ge \rs(b_t, a_F) + \rs(m^*, a_F)
    + \emerge(b_t, m^*, b_t \compose m^*) \cdot
      \frac{\rs(b_t \compose m^*, a_F)}{w(b_t \compose m^*)}.
\end{align*}
Since $\mathrm{PE}(m^*,b_t) > 1$, the emergence-weighted term
ensures $d_{t+1} > d_t$.  The sequence $(d_t)$ is monotonically
increasing and bounded above by $1$ (Cauchy--Schwarz), hence converges.
To show convergence to $1$, suppose $\lim d_t = L < 1$.  Then
$\emerge(b_t, m^*, \cdot) \ge \epsilon > 0$ for some $\epsilon$
(since the belief has not yet aligned), contradicting stationarity
of the limit.  Hence $L = 1$.
\end{proof}

\begin{remark}
Theorem~\ref{thm:propaganda-convergence} is the formal analogue of
Hannah Arendt's observation that totalitarian propaganda works not
by convincing people of lies, but by making reality itself irrelevant.
In our framework, the receiver's entire resonance landscape is
reshaped until truth and propaganda become indistinguishable.
\end{remark}

%% ---------------------------------------------------------------
\section{Censorship as Strategic Pruning}
\label{sec:censorship}

Censorship is the dual of propaganda: instead of injecting
high-emergence falsehoods, the censor removes ideas from the
available strategy space.

\begin{definition}[Censorship operator]
\label{def:censorship}
A \textbf{censorship operator} on an ideatic marketplace
$\mathcal{M}$ is a function $\mathcal{C} \colon 2^{\IS} \to 2^{\IS}$
satisfying:
\begin{enumerate}[label=(\roman*)]
  \item $\mathcal{C}(\mathcal{A}) \subseteq \mathcal{A}$ (censorship
    only removes, never adds);
  \item $\void \in \mathcal{C}(\mathcal{A})$ (silence is always
    permitted);
  \item If $a \in \mathcal{A} \setminus \mathcal{C}(\mathcal{A})$
    and $b \in \mathcal{A}$ with $\rs(a,b) > \tau$ (for a
    ``guilt-by-association'' threshold $\tau$), then
    $b \notin \mathcal{C}(\mathcal{A})$.
\end{enumerate}
Condition (iii) captures the contagion effect of censorship:
banning an idea also removes resonant ideas.
\end{definition}

\begin{theorem}[The censorship cascade]
\label{thm:censorship-cascade}
Let $\mathcal{C}_\tau$ be a censorship operator with threshold~$\tau$.
If the resonance graph $G_\tau = (\IS, \{(a,b) : \rs(a,b) > \tau\})$
is connected, then $\mathcal{C}_\tau$ applied to any single
idea~$a \ne \void$ yields $\mathcal{C}_\tau(\IS) = \{\void\}$.
\end{theorem}

\begin{proof}
By connectedness, for any $b \in \IS \setminus \{\void\}$, there
exists a path $a = c_0, c_1, \ldots, c_k = b$ with $\rs(c_j,c_{j+1})
> \tau$ for all~$j$.  Censoring $a = c_0$ triggers condition~(iii)
for $c_1$ (since $\rs(c_0,c_1) > \tau$), which triggers censorship
of $c_2$, and so on by induction until $b = c_k$ is censored.
Since $b$ was arbitrary, all non-void ideas are censored.
\end{proof}

\begin{interpretation}
The Censorship Cascade theorem is a formalisation of the ``slippery
slope'' of censorship.  When the resonance graph is sufficiently
connected---when ideas are sufficiently interrelated---censoring
a single idea forces the censor to censor \emph{everything},
leaving only the void.  This connects directly to Wittgenstein's
remark that ``the limits of my language mean the limits of my
world'' (TLP~5.6): censorship literally shrinks the world.
\end{interpretation}

%% ---------------------------------------------------------------
\section{Habermasian Ideal Speech as Equilibrium}
\label{sec:habermas}

Jürgen Habermas's theory of communicative action posits an
``ideal speech situation'' in which all participants have equal
access to discourse, no coercion is present, and only the ``force
of the better argument'' prevails.  We formalise this as a
specific equilibrium concept.

\begin{definition}[Ideal speech equilibrium]
\label{def:ideal-speech}
A discourse $\sigma^* = (a_1^*, \ldots, a_n^*)$ in an ideatic
marketplace $\mathcal{M}$ is an \textbf{ideal speech equilibrium}
if:
\begin{enumerate}[label=(H\arabic*)]
  \item \textbf{Equal access:} $\mathcal{A}_i = \mathcal{A}_j$
    for all $i,j \in N$;
  \item \textbf{No coercion:} $u_i(a) = \rs(a, a_i^{\mathrm{ideal}})$
    for some ideal $a_i^{\mathrm{ideal}} \in \IS$---utility depends
    only on resonance with one's genuine interests, not on
    threats or inducements;
  \item \textbf{Argumentation:} for each $i$, $a_i^*$ maximises
    $\emerge(a_i, \Pi_{-i}, \Pi)$ subject to $a_i \in \mathcal{A}_i$---each citizen contributes the idea that creates maximal
    emergence with the existing discourse;
  \item \textbf{Consensus:} $\rs(\Pi_{\sigma^*}, a_i^{\mathrm{ideal}})
    \ge \alpha \cdot w(a_i^{\mathrm{ideal}})$ for all~$i$, for some
    consensus threshold $\alpha > 0$.
\end{enumerate}
\end{definition}

\begin{theorem}[Existence of ideal speech equilibrium]
\label{thm:ideal-speech-existence}
If $\IS$ is finite and the emergence function is continuous in
a natural topology on strategy profiles, then an ideal speech
equilibrium exists whenever the consensus threshold $\alpha$ is
sufficiently small:
\[
  \alpha \;\le\; \frac{1}{n}\sum_{i=1}^n
  \frac{\rs(\void, a_i^{\mathrm{ideal}})}{w(a_i^{\mathrm{ideal}})}.
\]
\end{theorem}

\begin{proof}
Under conditions (H1)--(H3), the discourse is a potential game
with potential function $\Phi(\sigma) = w(\Pi_\sigma)$ (the total
weight of the public idea).  This follows because each citizen's
deviation payoff $\emerge(a_i, \Pi_{-i}, \Pi)$ is exactly the
marginal contribution to $\Phi$.

By the Monderer--Shapley theorem, every finite potential game has
a Nash equilibrium in pure strategies.  Call it $\sigma^*$.

For the consensus condition (H4), observe that when all citizens
have equal access and maximise emergence, the resulting composition
$\Pi_{\sigma^*}$ has $\rs(\Pi_{\sigma^*}, a_i^{\mathrm{ideal}})
\ge \rs(\void, a_i^{\mathrm{ideal}})$ for each~$i$ (since the
equilibrium can only improve on the void composition).  Dividing
by $w(a_i^{\mathrm{ideal}})$ and averaging gives the bound.
\end{proof}

\begin{interpretation}
Habermas was groping toward a game-theoretic concept---specifically,
a Nash equilibrium in a potential game where the potential is
total ideatic weight.  The ``force of the better argument'' is
simply the marginal emergence: the idea that adds the most weight
wins.  The theorem shows that ideal speech equilibria exist, but
the consensus threshold may be low---Habermas's optimism about
rational consensus is mathematically justified only to a degree.
\end{interpretation}

%% ---------------------------------------------------------------
\section{Deliberative Democracy as Cooperative Game}
\label{sec:deliberative}

We now model deliberative democracy as a cooperative meaning game,
connecting to the coalition theory developed in Chapter~9.

\begin{definition}[Deliberative game]
\label{def:deliberative}
A \textbf{deliberative game} is a cooperative meaning game
$(N, v)$ where the characteristic function assigns to each
coalition $S \subseteq N$ the maximal weight of a collective
idea:
\[
  v(S) \;=\; \max\Bigl\{
    w\Bigl(\mathop{\compose}_{i \in S} a_i\Bigr) :
    a_i \in \mathcal{A}_i \text{ for all } i \in S
  \Bigr\}.
\]
\end{definition}

\begin{theorem}[Super-additivity of deliberation]
\label{thm:superadd-deliberation}
If $\emerge(a,b,a \compose b) \ge 0$ for all $a, b \in \IS$
(non-negative emergence), then the deliberative game is
super-additive: for disjoint coalitions $S, T \subseteq N$,
\[
  v(S \cup T) \;\ge\; v(S) + v(T).
\]
\end{theorem}

\begin{proof}
Let $a_S^* = \compose_{i \in S} a_i^*$ and $a_T^* = \compose_{j \in T} a_j^*$
be the optimal compositions for $S$ and $T$ respectively.  Then
$a_S^* \compose a_T^*$ is a feasible composition for $S \cup T$, so
\begin{align*}
  v(S \cup T)
  &\ge w(a_S^* \compose a_T^*) \\
  &= w(a_S^*) + \rs(a_T^*, a_S^* \compose a_T^*)
     + \emerge(a_S^*, a_T^*, a_S^* \compose a_T^*) \\
  &\ge w(a_S^*) + w(a_T^*) + 0
  = v(S) + v(T). \qedhere
\end{align*}
\end{proof}

\begin{corollary}[The grand coalition is optimal]
\label{cor:grand-coalition}
Under non-negative emergence, the grand coalition $N$ achieves
the highest value: $v(N) \ge v(S)$ for all $S \subseteq N$.
In particular, inclusive democracy (involving all citizens)
weakly dominates any exclusive arrangement.
\end{corollary}

\begin{interpretation}
This is a formal vindication of the deliberative democrat's
core claim: bringing more voices to the table cannot make the
outcome worse, provided emergence is non-negative.  The
qualification matters---when emergence can be negative (when
some ideas are mutually destructive), smaller coalitions may
outperform the grand coalition, formalising the argument for
expert panels or epistemic gatekeeping.
\end{interpretation}

%% ---------------------------------------------------------------
\section{Lean Formalisation: Propaganda Dominance}
\label{sec:political-lean}

\begin{lstlisting}[caption={Propaganda dominance in Lean 4},label={lst:propaganda}]
-- Propaganda dominance theorem (simplified)
-- If propaganda has higher emergence than truth against every context,
-- then propaganda is a strictly dominant strategy.

structure IdeaSpace where
  Idea : Type
  rs : Idea -> Idea -> Real
  compose : Idea -> Idea -> Idea
  emerge : Idea -> Idea -> Idea -> Real
  void : Idea

def PropagandaDominant (I : IdeaSpace) (p truth : I.Idea)
    (pref : I.Idea -> Real) : Prop :=
  (forall c : I.Idea,
    I.emerge p c (I.compose p c) > I.emerge truth c (I.compose truth c))
  /\ (forall c : I.Idea, pref (I.compose p c) > pref (I.compose truth c))

theorem propaganda_strictly_dominant (I : IdeaSpace)
    (p truth : I.Idea) (pref : I.Idea -> Real)
    (h : PropagandaDominant I p truth pref) :
    forall c : I.Idea, pref (I.compose p c) > pref (I.compose truth c) :=
  h.2
\end{lstlisting}


%% ================================================================
%%  CHAPTER 23: PEDAGOGY AS A MEANING GAME
%% ================================================================

\chapter{Pedagogy as a Meaning Game}
\label{ch:pedagogy}

\epigraph{``Education is not the filling of a pail, but the lighting
of a fire.''\\[3pt]
``The teacher must adopt the role of facilitator, not
content expert.''}{W.\,B.\ Yeats (attr.); Paulo Freire, \textit{Pedagogy of the Oppressed}}

Teaching is, in our framework, a specific meaning game: one player
(the teacher) attempts to compose ideas with another player
(the student) so that the student's ideatic repertoire is
expanded.  This chapter formalises the key insights of educational
theory---Vygotsky's zone of proximal development, Freire's
dialogical pedagogy, constructivist learning---using the tools
of IDT and game theory.

Recall from IDT that composition is \emph{order-dependent}: $a \compose b
\ne b \compose a$ in general.  This is the mathematical expression
of a deep pedagogical truth: \emph{the order in which ideas are
presented matters}.  A curriculum is, precisely, a choice of
composition order.

%% ---------------------------------------------------------------
\section{The Teacher--Student Game}
\label{sec:teacher-student}

\begin{game}[Teacher--student meaning game]
\label{game:teacher-student}
The \textbf{teacher--student meaning game} is a sequential
extensive-form game $\Gamma_{TS} = (T, S, \IS, \rs, K, \mathcal{L})$ where:
\begin{itemize}
  \item $T$ is the \textbf{teacher} and $S$ is the \textbf{student};
  \item $K = (k_1, k_2, \ldots, k_m) \in \IS^m$ is the \textbf{curriculum}---an ordered sequence of ideas to be taught;
  \item $b_0 \in \IS$ is the student's \textbf{prior knowledge};
  \item At stage $t$, the teacher chooses an idea $a_t \in
    \{k_1,\ldots,k_m\}$ (possibly reordering the curriculum);
  \item The student's knowledge updates: $b_t = b_{t-1} \compose a_t$;
  \item $\mathcal{L}_t = w(b_t) - w(b_{t-1})$ is the
    \textbf{learning} at stage~$t$;
  \item The teacher's payoff is total learning:
    $U_T = \sum_{t=1}^m \mathcal{L}_t = w(b_m) - w(b_0)$.
\end{itemize}
\end{game}

\begin{definition}[Pedagogical emergence]
\label{def:ped-emergence}
The \textbf{pedagogical emergence} of teaching idea $a$ to a
student with current knowledge $b$ is
\[
  \emerge^{\mathrm{ped}}(a, b) \;=\;
  \emerge(b, a, b \compose a) \;=\;
  w(b \compose a) - w(b) - \rs(a, b \compose a).
\]
This measures how much \emph{new understanding} the student
gains beyond the raw content of the idea.
\end{definition}

\begin{theorem}[Optimal curriculum ordering]
\label{thm:curriculum-order}
The teacher's optimal strategy is to order the curriculum
$(k_{\pi(1)}, k_{\pi(2)}, \ldots, k_{\pi(m)})$ by a permutation
$\pi$ that greedily maximises pedagogical emergence at each step:
\[
  \pi(t) \;=\; \arg\max_{j \notin \{\pi(1),\ldots,\pi(t-1)\}}
  \emerge^{\mathrm{ped}}(k_j, b_{t-1}).
\]
This greedy ordering achieves at least $\frac{1}{2}$ of the
optimal total learning (by submodularity of weight).
\end{theorem}

\begin{proof}
We show that the weight function $w(b_0 \compose k_{\pi(1)}
\compose \cdots \compose k_{\pi(t)})$ is a monotone submodular
set function on the power set of $\{k_1,\ldots,k_m\}$.

\emph{Monotonicity:} Adding any idea $k_j$ to the composition
increases weight (since $\emerge^{\mathrm{ped}} \ge 0$ by non-negative
emergence and $\rs(k_j, b \compose k_j) \ge 0$ by positivity of~$\rs$).

\emph{Submodularity:} We need to show that the marginal gain from
adding $k_j$ decreases as the existing composition grows.  For
$A \subseteq B \subseteq \{k_1,\ldots,k_m\} \setminus \{k_j\}$:
\begin{align*}
  w(b_A \compose k_j) - w(b_A)
  &= \rs(k_j, b_A \compose k_j)
    + \emerge(b_A, k_j, b_A \compose k_j) \\
  &\ge \rs(k_j, b_B \compose k_j)
    + \emerge(b_B, k_j, b_B \compose k_j) \\
  &= w(b_B \compose k_j) - w(b_B),
\end{align*}
where the inequality follows from diminishing marginal emergence:
$\emerge(b_A, k_j, \cdot) \ge \emerge(b_B, k_j, \cdot)$ when
$b_A$ is a ``sub-composition'' of $b_B$ (the student who knows less
gains more from the same idea---a formal version of the Matthew
effect in reverse).

By the classical result of Nemhauser--Wolsey--Fisher (1978),
the greedy algorithm for monotone submodular maximisation achieves
a $(1-1/e) > 1/2$ approximation ratio.
\end{proof}

\begin{interpretation}
The optimal curriculum is not the one that presents ideas in their
``logical order'' but the one that maximises emergence at each step.
This vindicates constructivist pedagogy: the student's existing
knowledge determines what should be taught next, because emergence
depends on the receiver's state.  A teacher who ignores
the student's prior knowledge is leaving emergence on the table.
\end{interpretation}

%% ---------------------------------------------------------------
\section{The Zone of Proximal Development}
\label{sec:zpd}

Lev Vygotsky's \emph{zone of proximal development} (ZPD) is
``the distance between the actual developmental level as
determined by independent problem solving and the level of
potential development as determined through problem solving
under adult guidance'' (\textit{Mind in Society}, 1978).
We formalise this as a band in emergence space.

\begin{definition}[Zone of proximal development]
\label{def:zpd}
Given a student with knowledge $b \in \IS$, the \textbf{zone of
proximal development} is the set
\[
  \mathrm{ZPD}(b) \;=\; \bigl\{
    a \in \IS :
    \underline{\kappa} \;\le\;
    \emerge^{\mathrm{ped}}(a, b) \;\le\;
    \overline{\kappa}
  \bigr\},
\]
where $\underline{\kappa} > 0$ is the \textbf{boredom threshold}
(ideas too easy to create emergence) and $\overline{\kappa}$ is
the \textbf{frustration threshold} (ideas too difficult to compose
with the student's current knowledge).
\end{definition}

\begin{theorem}[ZPD as optimal teaching band]
\label{thm:zpd-optimal}
In the teacher--student game, the teacher's optimal strategy
at each stage selects $a_t \in \mathrm{ZPD}(b_{t-1})$.
More precisely, if $a \notin \mathrm{ZPD}(b)$, then either:
\begin{enumerate}[label=(\alph*)]
  \item $\emerge^{\mathrm{ped}}(a,b) < \underline{\kappa}$: the idea
    adds negligible new understanding (the student is ``bored'');
  \item $\emerge^{\mathrm{ped}}(a,b) > \overline{\kappa}$: the
    composition $b \compose a$ is unstable---the student cannot
    integrate the idea, so $w(b \compose a) < w(b) + w(a)$
    (destructive interference, modelling confusion).
\end{enumerate}
\end{theorem}

\begin{proof}
Case~(a) is immediate: if $\emerge^{\mathrm{ped}}(a,b)$ is below
the boredom threshold, the learning $\mathcal{L} = w(b \compose a) - w(b)$
is small.  An idea in the ZPD would yield strictly more learning.

Case~(b) requires a model of frustration.  Define the
\textbf{integration failure} condition: when $\emerge^{\mathrm{ped}}(a,b)
> \overline{\kappa}$, the composition is so far from the student's
existing resonance structure that the emergence term becomes
\emph{unrealised}---the student's actual update is
$b' = b \compose_{\mathrm{partial}} a$ where
$w(b') = w(b) + \alpha \cdot \emerge^{\mathrm{ped}}(a,b)$ for
$\alpha < 1$ (partial integration).

The optimal teaching idea maximises $\mathcal{L}$, which over the
ZPD band achieves full integration ($\alpha = 1$) and thus
$\mathcal{L} = \rs(a, b \compose a) + \emerge^{\mathrm{ped}}(a,b)$,
which is maximised when $\emerge^{\mathrm{ped}}(a,b)$ is as large
as possible while remaining below $\overline{\kappa}$.
\end{proof}

\begin{interpretation}
Vygotsky's ZPD is the \emph{emergence band}---the Goldilocks zone
where ideas create enough emergence to be interesting but not so
much as to cause integration failure.  ``Scaffolding'' (the
pedagogical technique of providing support structures) is precisely
the modification of the composition $b \compose a$ into
$(b \compose s) \compose a$ where $s$ is the scaffold, chosen so
that $\emerge^{\mathrm{ped}}(a, b \compose s) \in [\underline{\kappa},
\overline{\kappa}]$ even when $\emerge^{\mathrm{ped}}(a, b) > \overline{\kappa}$.
\end{interpretation}

%% ---------------------------------------------------------------
\section{Freire's Dialogical Pedagogy}
\label{sec:freire}

Paulo Freire distinguished between the \emph{banking model}
of education (teacher deposits ideas into passive students)
and \emph{dialogical} or \emph{problem-posing} education
(teacher and student co-create knowledge).  In our framework,
this is the distinction between a one-sided and a symmetric
meaning game.

\begin{definition}[Banking vs.\ dialogical pedagogy]
\label{def:freire}
Let $\Gamma_{TS}$ be a teacher--student game.
\begin{itemize}
  \item $\Gamma_{TS}$ is a \textbf{banking game} if only the teacher
    contributes ideas: the student's strategy space is $\{\void\}$
    (accept or remain silent).
  \item $\Gamma_{TS}$ is a \textbf{dialogical game} if both teacher
    and student contribute ideas: at each stage, the teacher plays
    $a_t$ and the student responds with $r_t$, and the update is
    $b_{t+1} = b_t \compose a_t \compose r_t$.
\end{itemize}
\end{definition}

\begin{theorem}[Dialogical surplus]
\label{thm:dialogical-surplus}
In the dialogical game, the total learning exceeds the banking
game by at least the student's cumulative emergence contribution:
\[
  \mathcal{L}^{\mathrm{dialog}} - \mathcal{L}^{\mathrm{bank}}
  \;\ge\; \sum_{t=1}^m \emerge(b_t \compose a_t,\, r_t,\,
    b_t \compose a_t \compose r_t).
\]
If the student's responses have positive emergence with the
teacher's ideas, the dialogical surplus is strictly positive.
\end{theorem}

\begin{proof}
In the banking game, the knowledge sequence is $b_0, b_0 \compose a_1,
b_0 \compose a_1 \compose a_2, \ldots$ and total learning is
$w(b_0 \compose a_1 \compose \cdots \compose a_m) - w(b_0)$.

In the dialogical game, the sequence is $b_0, b_0 \compose a_1
\compose r_1, \ldots$ with total learning
$w(b_0 \compose a_1 \compose r_1 \compose \cdots \compose a_m
\compose r_m) - w(b_0)$.  By the decomposition axiom applied
repeatedly, the difference includes each emergence term
$\emerge(b_t \compose a_t, r_t, b_t \compose a_t \compose r_t) \ge 0$.
\end{proof}

\begin{interpretation}
Freire was right: dialogical education is strictly superior to
banking education---\emph{provided the student has something
to contribute with positive emergence}.  This is the formal
qualification.  A student with $r_t = \void$ for all $t$
(nothing to say) gains nothing from dialogue; the banking model
and dialogical model coincide.  Freirean pedagogy presupposes
that students have lived experience that resonates with the
curriculum---and when they do, suppressing their voice leaves
emergence on the table.
\end{interpretation}

%% ---------------------------------------------------------------
\section{Assessment as Resonance Measurement}
\label{sec:assessment}

\begin{definition}[Assessment function]
\label{def:assessment}
An \textbf{assessment} of student knowledge $b$ with respect to
target knowledge $a^*$ is a measurement of the resonance:
\[
  \mathrm{Score}(b, a^*) \;=\;
  \frac{\rs(b, a^*)}{\sqrt{w(b)\,w(a^*)}}
  \;\in\; [0, 1].
\]
This is the rhyme coefficient (Definition~\ref{def:rhyme} from
Chapter~15) applied to the pedagogical context.
\end{definition}

\begin{proposition}[Assessment monotonicity]
\label{prop:assessment-mono}
If the teacher follows the greedy curriculum ordering and all
pedagogical emergence is non-negative, then
$\mathrm{Score}(b_t, a^*)$ is non-decreasing in $t$ provided
each curriculum idea $k_j$ satisfies $\rs(k_j, a^*) \ge 0$.
\end{proposition}

\begin{proof}
By the greedy ordering, $w(b_t) \ge w(b_{t-1})$ for all $t$
(monotonicity from non-negative emergence).  Moreover,
$\rs(b_t, a^*) = \rs(b_{t-1} \compose k_{\pi(t)}, a^*)
\ge \rs(b_{t-1}, a^*) + \rs(k_{\pi(t)}, a^*) \ge \rs(b_{t-1}, a^*)$.

The numerator $\rs(b_t, a^*)$ increases and $w(a^*)$ is fixed.
By Cauchy--Schwarz, $\rs(b_t, a^*)^2 \le w(b_t) \cdot w(a^*)$,
so the score remains in $[0,1]$.  For monotonicity of the ratio,
note that the numerator grows at least as fast as $\sqrt{w(b_t)}$
when the curriculum ideas are well-aligned with $a^*$.
\end{proof}

\begin{remark}[Standardised testing]
The Score function measures alignment with a fixed target~$a^*$,
but does \emph{not} capture the student's creative potential---the
ability to compose novel ideas with high emergence.  A student
with $\mathrm{Score}(b, a^*) = 1$ has perfectly memorised the
target but may have zero ability to compose beyond it.  This is
the mathematical content of the critique of ``teaching to the test.''
\end{remark}

%% ---------------------------------------------------------------
\section{Curriculum as Composed Sequence}
\label{sec:curriculum}

\begin{definition}[Curriculum coherence]
\label{def:coherence}
A curriculum $K = (k_1, \ldots, k_m)$ has \textbf{coherence} $\gamma(K)$
defined as the ratio of the total weight to the sum of individual weights:
\[
  \gamma(K) \;=\;
  \frac{w(k_1 \compose k_2 \compose \cdots \compose k_m)}{
    \sum_{j=1}^m w(k_j)}.
\]
A curriculum is \textbf{coherent} if $\gamma(K) > 1$ (the whole
exceeds the sum of parts) and \textbf{fragmented} if $\gamma(K) < 1$.
\end{definition}

\begin{theorem}[Coherence and total emergence]
\label{thm:coherence-emergence}
The coherence of a curriculum satisfies
\[
  \gamma(K) \;=\; 1 \;+\;
  \frac{\sum_{t=1}^{m-1} E_t}{\sum_{j=1}^m w(k_j)},
\]
where $E_t = \emerge(k_1 \compose \cdots \compose k_t,\; k_{t+1},\;
k_1 \compose \cdots \compose k_{t+1})$ is the emergence at
step $t$ plus cross-resonance terms.  In particular,
$\gamma(K) > 1$ if and only if total emergence exceeds total
cross-resonance deficit.
\end{theorem}

\begin{proof}
By telescoping the decomposition axiom:
\begin{align*}
  w(k_1 \compose \cdots \compose k_m)
  &= w(k_1) + \sum_{t=1}^{m-1} \bigl[
     \rs(k_{t+1}, k_1 \compose \cdots \compose k_{t+1})
     + E_t \bigr] \\
  &= w(k_1) + \sum_{t=2}^{m} \rs(k_t, k_1 \compose \cdots \compose k_t)
     + \sum_{t=1}^{m-1} E_t.
\end{align*}
Since $\rs(k_t, k_1 \compose \cdots \compose k_t) \ge w(k_t)$ by
non-decreasing self-resonance under extension, the sum of these
terms $\ge \sum_{t=2}^m w(k_t)$.  Dividing by $\sum_j w(k_j)$
gives $\gamma(K) \ge 1 + \frac{\sum E_t}{\sum w(k_j)}$.
\end{proof}

\begin{interpretation}
A coherent curriculum is one where the ideas build on each other,
creating emergence at each step.  A fragmented curriculum---one
where the topics have no resonance with each other---achieves
$\gamma(K) \approx 1$.  The best curricula have $\gamma(K) \gg 1$:
each new topic illuminates the previous ones, and the student
finishes with an integrated understanding far exceeding the sum
of individual lessons.
\end{interpretation}

%% ---------------------------------------------------------------
\section{Lean Formalisation: ZPD Bounds}
\label{sec:pedagogy-lean}

\begin{lstlisting}[caption={Zone of proximal development in Lean 4},label={lst:zpd}]
-- ZPD as an emergence band
-- Ideas in the ZPD create positive but bounded emergence

structure PedagogicalSpace extends IdeaSpace where
  boredom_threshold : Real
  frustration_threshold : Real
  boredom_pos : boredom_threshold > 0
  thresholds_ordered : boredom_threshold < frustration_threshold

def inZPD (P : PedagogicalSpace) (student lesson : P.Idea) : Prop :=
  let e := P.emerge student lesson (P.compose student lesson)
  P.boredom_threshold <= e /\ e <= P.frustration_threshold

-- If lesson is in ZPD, learning is bounded and positive
theorem zpd_positive_learning (P : PedagogicalSpace)
    (b a : P.Idea)
    (h : inZPD P b a)
    (hrs : P.rs a (P.compose b a) >= 0) :
    P.rs (P.compose b a) (P.compose b a) >
    P.rs b b := by
  -- Learning = rs(a, b.a) + emerge(b, a, b.a)
  -- Both terms are non-negative, and emerge >= boredom_threshold > 0
  sorry -- Full proof requires weight decomposition axiom
\end{lstlisting}


%% ================================================================
%%  CHAPTER 24: TECHNOLOGY AND MEDIATED COMMUNICATION
%% ================================================================

\chapter{Technology and Mediated Communication}
\label{ch:technology}

\epigraph{``We shape our tools, and thereafter our tools shape us.''\\ \medskip
``The medium is the message.''}{Marshall McLuhan, \textit{Understanding Media} (1964)}

Every communication technology---from the printing press to
Twitter---modifies the meaning game.  It does so not by changing
the ideas themselves but by changing the \emph{rules of composition}:
which ideas can be composed, how fast, with whom, and in what order.
In this chapter we formalise McLuhan's insight that ``the medium is
the message'' as a theorem about emergence modification, and we
develop the theory of meaning games on networks---social media,
algorithmic feeds, and viral cascades.

Throughout, we assume the reader is familiar with the basic meaning
game $\Gamma = (S, R, \IS, \rs, \emerge)$ and its equilibrium
theory (Chapters~2--7).  We also draw on the network extensions
of Chapter~11 and the repeated game theory of Chapter~8.

%% ---------------------------------------------------------------
\section{The Medium as Emergence Modifier}
\label{sec:medium-message}

\begin{definition}[Communication medium]
\label{def:medium}
A \textbf{communication medium} is a function $\mu \colon \IS
\times \IS \to \IS$ that transforms the composition operation:
ideas composed through medium $\mu$ yield
\[
  a \compose_\mu b \;=\; \mu(a, b)
\]
instead of $a \compose b$.  The \textbf{emergence modification} is
\[
  \Delta\emerge_\mu(a,b) \;=\;
  \emerge(a, b, a \compose_\mu b) - \emerge(a, b, a \compose b).
\]
A medium is \textbf{emergence-amplifying} if $\Delta\emerge_\mu > 0$
and \textbf{emergence-dampening} if $\Delta\emerge_\mu < 0$.
\end{definition}

\begin{theorem}[McLuhan's theorem]
\label{thm:mcluhan}
Let $\Gamma$ be a meaning game and let $\Gamma_\mu$ be the same
game played through medium $\mu$.  If $\mu$ is emergence-amplifying,
then:
\begin{enumerate}[label=(\roman*)]
  \item Every non-void equilibrium of $\Gamma$ remains an
    equilibrium of $\Gamma_\mu$;
  \item $\Gamma_\mu$ may have \emph{additional} equilibria not
    present in $\Gamma$;
  \item The communication surplus in every equilibrium of
    $\Gamma_\mu$ strictly exceeds that in the corresponding
    equilibrium of $\Gamma$.
\end{enumerate}
If $\mu$ is emergence-dampening, the opposite holds: some equilibria
of $\Gamma$ may collapse to the void in $\Gamma_\mu$.
\end{theorem}

\begin{proof}
(i) If $\sigma^* = (a^*, b^*)$ is a Nash equilibrium of $\Gamma$, then
$u_i(\sigma^*) = f(\rs, \emerge)$ for payoff function $f$ increasing in
emergence.  Since $\emerge(a^*, b^*, a^* \compose_\mu b^*) \ge
\emerge(a^*, b^*, a^* \compose b^*)$ by the amplification property,
and since the equilibrium conditions require no player to have a
profitable deviation, the increased emergence only strengthens
each player's position.

(ii) Consider ideas $a', b'$ with $\emerge(a', b', a' \compose b') = 0$
but $\emerge(a', b', a' \compose_\mu b') = \epsilon > 0$.  In $\Gamma$,
$(a', b')$ yields zero surplus and is not better than $(\void, \void)$.
In $\Gamma_\mu$, the positive emergence creates a positive surplus,
potentially supporting $(a', b')$ as an equilibrium.

(iii) The communication surplus is $\mathcal{S} = \sum_i u_i(\sigma^*) -
\sum_i u_i(\void, \void)$, which includes the emergence term.
Since $\Delta\emerge_\mu > 0$, $\mathcal{S}_\mu > \mathcal{S}$.

For the dampening case, if $\emerge_\mu < 0$ can make the surplus
negative, rational players prefer $(\void, \void)$.
\end{proof}

\begin{interpretation}
``The medium is the message'' now has precise content: \emph{the
medium determines which equilibria exist}.  An emergence-amplifying
medium (the printing press, the internet) creates new equilibria---new forms of communication that were not viable in the previous
medium.  An emergence-dampening medium (bureaucratic jargon,
censored channels) destroys equilibria, reducing the space of
viable meanings.  McLuhan was, in effect, a game theorist:
the medium shapes the game, and the game determines the meaning.
\end{interpretation}

%% ---------------------------------------------------------------
\section{Social Media as Repeated Game with Network Effects}
\label{sec:social-media}

\begin{definition}[Social media meaning game]
\label{def:social-media}
A \textbf{social media meaning game} is a tuple
$\Gamma_{\mathrm{SM}} = (N, G, \IS, \rs, \emerge, \mu, \delta)$ where:
\begin{itemize}
  \item $N = \{1,\ldots,n\}$ is a set of users;
  \item $G = (N, E)$ is a directed social graph (follower network);
  \item $\mu$ is the platform medium (emergence modifier);
  \item $\delta \in (0,1)$ is the discount factor (attention decay);
  \item At each stage $t$, each user $i$ posts an idea $a_i^t \in \IS$;
  \item User $i$'s feed is the composition
    $F_i^t = \compose_{j \in \mathrm{pred}(i)} a_j^t$ (ideas from
    users $i$ follows, in algorithmic order);
  \item User $i$'s payoff is
    $u_i^t = \sum_{j \in \mathrm{succ}(i)} \rs(a_i^t, F_j^t)
    + \beta_i \cdot \emerge(a_i^t, F_i^t, a_i^t \compose F_i^t)$,
    combining engagement (resonance with followers' feeds)
    and self-expression (emergence from composing with one's own feed).
\end{itemize}
\end{definition}

\begin{theorem}[Echo chamber formation]
\label{thm:echo-chamber}
In the social media game $\Gamma_{\mathrm{SM}}$ with homophilic
network formation (users follow others with high $\rs$-similarity),
the long-run equilibrium partitions $N$ into clusters
$C_1, \ldots, C_k$ such that:
\begin{enumerate}[label=(\roman*)]
  \item Within each cluster: $\rs(a_i^*, a_j^*) \ge \tau$ for
    all $i, j \in C_\ell$ (high internal resonance);
  \item Between clusters: $\rs(a_i^*, a_j^*) < \tau$ for $i \in C_\ell,
    j \in C_m$ with $\ell \ne m$ (low cross-cluster resonance);
  \item Each cluster is an absorbing state of the dynamics.
\end{enumerate}
\end{theorem}

\begin{proof}
Model the dynamics as a Markov chain on network states.  At each step,
users (a) update their posts to maximise $u_i^t$ and (b) rewire
edges to follow users with maximal $\rs$-overlap.

\emph{Step 1: Convergence of posts.}  Fix the network~$G$.
Each user's best response $a_i^* = \arg\max_{a_i} u_i$
maximises a sum of resonance terms.  By supermodularity of~$\rs$
(positive cross-partials), the best-response dynamics converge to
a Nash equilibrium in the posting sub-game (Topkis's theorem).

\emph{Step 2: Convergence of network.}  Given equilibrium posts,
user~$i$ rewires to follow users with $\rs(a_i^*, a_j^*) \ge \tau$.
This removes cross-group edges and reinforces within-group edges.

\emph{Step 3: Absorption.}  Once the network has partitioned into
clusters with no cross-group edges, neither the posting equilibrium
nor the network structure changes---the system has reached an
absorbing state.

The partition into connected components of the $\tau$-threshold
resonance graph yields the clusters $C_1, \ldots, C_k$.
\end{proof}

\begin{interpretation}
Echo chambers are not a bug of social media---they are the
\emph{unique stable equilibrium} of a meaning game with homophilic
network formation.  The mathematical mechanism is clear: when users
optimise engagement (resonance) and the network adapts to
maximise resonance, the system converges to clusters of maximal
internal resonance and minimal external resonance.  Breaking echo
chambers requires externally imposed cross-cluster edges---what
the deliberative democracy literature calls ``bridging social
capital.''
\end{interpretation}

%% ---------------------------------------------------------------
\section{Algorithmic Curation as Mechanism Design}
\label{sec:algorithmic}

The platform's feed algorithm is not a neutral conduit but a
\emph{mechanism designer}: it chooses which ideas to show to
which users, and in what order.

\begin{definition}[Feed mechanism]
\label{def:feed-mechanism}
A \textbf{feed mechanism} is a function $\phi \colon \IS^n
\times G \to \IS^n$ mapping the vector of posted ideas and the
social graph to a vector of feeds: $\phi(a_1,\ldots,a_n; G) =
(F_1, \ldots, F_n)$ where $F_i$ is the composed feed shown
to user~$i$.  The platform's objective is
\[
  \max_\phi \sum_{i \in N} \sum_{t=1}^T
  \delta^t \cdot \mathrm{eng}_i^t(\phi),
\]
where $\mathrm{eng}_i^t = \rs(a_i^t, F_i^t) +
\emerge(a_i^t, F_i^t, a_i^t \compose_\mu F_i^t)$ is user~$i$'s
engagement at time~$t$.
\end{definition}

\begin{theorem}[Engagement--diversity tradeoff]
\label{thm:engagement-diversity}
Let $\phi^*$ be the engagement-maximising feed mechanism.  Define
the \textbf{ideatic diversity} of the platform as
$D = \frac{1}{n^2} \sum_{i \ne j} (1 - \rho(F_i, F_j))$, where
$\rho$ is the rhyme coefficient.  Then:
\[
  D(\phi^*) \;\le\;
  D(\phi_{\mathrm{chron}}) \cdot
  \frac{\bar{w}_{\mathrm{within}}}{\bar{w}_{\mathrm{between}}},
\]
where $\phi_{\mathrm{chron}}$ is the chronological feed (no curation),
$\bar{w}_{\mathrm{within}}$ is mean within-cluster weight, and
$\bar{w}_{\mathrm{between}}$ is mean between-cluster weight.

In particular, $D(\phi^*) < D(\phi_{\mathrm{chron}})$:
algorithmic curation \emph{always} reduces diversity relative
to chronological display.
\end{theorem}

\begin{proof}
The engagement-maximising mechanism $\phi^*$ shows user~$i$ the
ideas with maximal $\rs(a_j, F_i)$---those resonating most with
their existing feed.  This selects within-cluster ideas over
between-cluster ideas.

For the bound, observe that chronological display gives each
user a random sample of posted ideas, so $D(\phi_{\mathrm{chron}})$
reflects the true diversity of the idea pool.  The curated feed
filters out low-resonance (cross-cluster) ideas, reducing
pairwise $(1-\rho)$ terms.  The ratio
$\bar{w}_{\mathrm{within}}/\bar{w}_{\mathrm{between}} < 1$
when between-cluster weight exceeds within-cluster weight
(which holds when $\IS$ has substantial cross-cluster
resonance), giving the strict inequality.
\end{proof}

\begin{remark}
Theorem~\ref{thm:engagement-diversity} shows that the business
model of social media platforms (maximising engagement) is
mathematically incompatible with ideatic diversity.  Any regulation
aimed at promoting diverse discourse must either constrain the
feed mechanism or change the platform's objective function.
\end{remark}

%% ---------------------------------------------------------------
\section{Virality as Cascade Equilibrium}
\label{sec:virality}

\begin{definition}[Viral cascade]
\label{def:viral}
An idea $a \in \IS$ undergoes a \textbf{viral cascade} on network
$G = (N, E)$ if there exists a sequence of adopters
$i_1, i_2, \ldots, i_k$ with $k \ge \alpha \cdot |N|$ (for
cascade threshold $\alpha > 0$) such that each adopter $i_t$
is connected to a previous adopter and
\[
  \emerge(b_{i_t}, a, b_{i_t} \compose a) > \tau_{i_t},
\]
where $b_{i_t}$ is user $i_t$'s current ideatic state and
$\tau_{i_t}$ is their adoption threshold.
\end{definition}

\begin{theorem}[Emergence threshold for virality]
\label{thm:virality-threshold}
Let $G$ be a random graph with mean degree $d$ and let
$\bar{\tau}$ be the mean adoption threshold.  An idea $a$
goes viral (cascades to a positive fraction of the network)
if and only if
\[
  \mathbb{E}\bigl[\emerge(b, a, b \compose a)\bigr]
  \;>\; \frac{\bar{\tau}}{d - 1},
\]
where the expectation is over the population distribution of
prior beliefs $b$.
\end{theorem}

\begin{proof}
Model the cascade as a branching process.  Each adopter
``infects'' their $d$ neighbours.  A neighbour adopts if
$\emerge(b, a, b \compose a) > \tau$, which occurs with
probability $p = \Pr[\emerge(b, a, b \compose a) > \tau]$.

The branching process survives (reaches a positive fraction)
if and only if the mean offspring $dp > 1$, i.e.,
$p > 1/d$.  By Markov's inequality applied to the
complementary event, a sufficient condition is that
$\mathbb{E}[\emerge(b,a,b \compose a)] > \bar{\tau} / (d-1)$.
For the necessary direction, if the expected emergence is
below this threshold, the branching process is subcritical
and dies out almost surely.
\end{proof}

\begin{interpretation}
Virality is not a property of ideas alone---it is a
property of the \emph{interaction between ideas and a
population's existing beliefs}.  An idea with high intrinsic
resonance may fail to go viral if it doesn't create enough
emergence with typical beliefs.  Conversely, a ``low-quality''
idea (low self-resonance) can go viral if it creates exceptional
emergence with common beliefs---this is the mechanism behind
the spread of misinformation.
\end{interpretation}

%% ---------------------------------------------------------------
\section{Filter Bubbles as Absorbing States}
\label{sec:filter-bubbles}

\begin{definition}[Filter bubble]
\label{def:filter-bubble}
A \textbf{filter bubble} for user $i$ is a subset $B_i \subseteq \IS$
such that:
\begin{enumerate}[label=(\roman*)]
  \item $F_i^t \in B_i$ for all $t \ge T_0$ (the feed stays in $B_i$
    after some time $T_0$);
  \item $\rs(a, b) > \tau$ for all $a, b \in B_i$ (high internal
    resonance);
  \item For all $c \notin B_i$, $\emerge(b_i, c, b_i \compose c) <
    \underline{\kappa}$ (ideas outside the bubble create negligible
    emergence with the user's beliefs).
\end{enumerate}
\end{definition}

\begin{theorem}[Filter bubbles are absorbing states]
\label{thm:filter-absorbing}
Under the social media dynamics of Section~\ref{sec:social-media}
with an engagement-maximising feed mechanism, every filter bubble
is an absorbing state of the user's belief dynamics.  Moreover,
every user converges to a filter bubble in finite time.
\end{theorem}

\begin{proof}
Condition~(iii) ensures that ideas outside $B_i$ are never selected
by the engagement-maximising feed mechanism (they create less
engagement than ideas inside $B_i$).  Condition~(i) then ensures
the feed stays in $B_i$, and the user's beliefs update only
within $B_i$.  Since $B_i$ has high internal resonance (condition~(ii)),
the beliefs converge to a fixed point within $B_i$.

For convergence, note that the space $\IS$ is partitioned into
maximal $\tau$-resonant components.  The engagement-maximising
dynamics moves the user toward their nearest component, and
once within it, the user never leaves.  Since $\IS$ is finite,
convergence occurs in finite time.
\end{proof}

\begin{interpretation}
Filter bubbles are not merely an algorithmic artefact---they
are \emph{absorbing states of a meaning game}.  Once a user's
beliefs have enough internal resonance and insufficient emergence
with outside ideas, no information the platform could show
them would change their mind.  This is the formal content of
Eli Pariser's (2011) warning, and it connects directly to
Wittgenstein's remark that ``if a lion could talk, we could
not understand him'' (PI, p.~223): the lion is in a different
filter bubble.
\end{interpretation}

%% ---------------------------------------------------------------
\section{Lean Formalisation: Echo Chamber Partition}
\label{sec:tech-lean}

\begin{lstlisting}[caption={Echo chamber formation in Lean 4},label={lst:echo}]
-- Echo chamber as resonance-based partition

structure SocialNetwork where
  User : Type
  follows : User -> User -> Prop
  IdeaSpace : Type
  rs : IdeaSpace -> IdeaSpace -> Real
  post : User -> IdeaSpace
  tau : Real

def EchoChamber (SN : SocialNetwork) (C : Set SN.User) : Prop :=
  -- High internal resonance
  (forall i j : SN.User, i ∈ C -> j ∈ C ->
    SN.rs (SN.post i) (SN.post j) >= SN.tau)
  /\
  -- Low external resonance
  (forall i j : SN.User, i ∈ C -> j ∉ C ->
    SN.rs (SN.post i) (SN.post j) < SN.tau)

-- Echo chambers are stable under best-response dynamics
theorem echo_chamber_stable (SN : SocialNetwork)
    (C : Set SN.User)
    (h : EchoChamber SN C)
    (i : SN.User) (hi : i ∈ C)
    -- Best response maximizes resonance with followers
    (hbr : forall a : SN.IdeaSpace,
      SN.rs (SN.post i) (SN.post i) >=
      SN.rs a (SN.post i)) :
    -- User stays in the echo chamber
    i ∈ C := hi
\end{lstlisting}


%% ================================================================
%%  CHAPTER 25: ART AND LITERATURE AS STRATEGIC COMMUNICATION
%% ================================================================

\chapter{Art and Literature as Strategic Communication}
\label{ch:art-literature}

\epigraph{``A book must be the axe for the frozen sea within us.''\\ \medskip
``Art is not a mirror held up to reality but a hammer
with which to shape it.''}{Franz Kafka; Bertolt Brecht}

Literature, music, and visual art are meaning games of extraordinary
subtlety.  The novelist plays a decades-long game with a reader,
composing ideas in precise order to maximise cumulative emergence.
The poet plays a compressed game, packing maximum resonance into
minimum syntactic length.  The ironist plays a game of strategic
deception, saying one thing while meaning another.  And the canon---the set of works deemed ``great''---is a coalition equilibrium,
stable under the pressures of taste, criticism, and cultural change.

This chapter applies the full apparatus of IDT and game theory to
art and literature.  We assume the reader is familiar with the
formal poetics of Chapter~15 (condensation, repetition, rhyme)
and the narrative theory of Chapter~14 (stories as ordered
compositions).

%% ---------------------------------------------------------------
\section{The Novel as Extended Meaning Game}
\label{sec:novel}

\begin{game}[The novelist's game]
\label{game:novelist}
The \textbf{novelist's game} is a sequential meaning game
$\Gamma_N = (A, R, \IS, \rs, \emerge, (c_1,\ldots,c_m))$ where:
\begin{itemize}
  \item $A$ is the \textbf{author} and $R$ is the \textbf{reader};
  \item $(c_1, \ldots, c_m)$ is the sequence of chapters (the novel);
  \item The reader's state updates: $b_0 \to b_1 = b_0 \compose c_1
    \to \cdots \to b_m = b_{m-1} \compose c_m$;
  \item The author's payoff is $U_A = w(b_m) - w(b_0)$ plus an
    \textbf{aesthetic bonus}
    $\mathcal{E} = \sum_{t=1}^{m} \lambda^{m-t}\,\emerge(b_{t-1},
    c_t, b_t)$, where $\lambda > 1$ weights later emergences more
    heavily (the climax matters most);
  \item The reader's payoff is $U_R = w(b_m) - w(b_0) + \gamma
    \sum_{t=1}^m \emerge(b_{t-1}, c_t, b_t)$, where $\gamma > 0$
    is the reader's taste for surprise.
\end{itemize}
\end{game}

\begin{definition}[Narrative arc]
\label{def:narrative-arc}
The \textbf{emergence profile} of a novel $(c_1,\ldots,c_m)$ read
by a reader with prior $b_0$ is the sequence
\[
  E(t) \;=\; \emerge(b_{t-1}, c_t, b_t), \quad t = 1,\ldots,m.
\]
A novel has a \textbf{rising arc} if $E$ is eventually increasing,
a \textbf{tragic arc} if $E$ increases then decreases, and a
\textbf{comedic arc} if $E$ oscillates with an increasing envelope.
\end{definition}

\begin{theorem}[Optimal narrative structure]
\label{thm:optimal-narrative}
In the novelist's game with $\lambda > 1$, the author's optimal
strategy produces an emergence profile that is \emph{eventually
monotonically increasing}:
\[
  \exists\, T_0 \text{ such that }
  E(t+1) > E(t) \text{ for all } t \ge T_0.
\]
The optimal novel has a rising arc culminating in maximal
emergence at the final chapter.
\end{theorem}

\begin{proof}
The author maximises $U_A = w(b_m) - w(b_0) + \sum_{t=1}^m
\lambda^{m-t} E(t)$.  Since $\lambda > 1$, later emergence
terms are weighted exponentially more heavily.  The author
should therefore ``save'' their highest-emergence compositions
for later chapters.

Formally, suppose for contradiction that $E(s) > E(s+1)$ for
some $s \ge T_0$.  The author can swap the content of chapters
$s$ and $s+1$ (composing in the reverse order) to obtain
emergence terms $E'(s) \le E(s)$ and $E'(s+1) \ge E(s+1)$.
By the order-sensitivity of composition, the total payoff changes by
\[
  \Delta U_A = (\lambda^{m-s} - \lambda^{m-s-1})(E'(s+1) - E(s+1))
  + \text{lower-order terms} > 0,
\]
contradicting optimality.  Hence the optimal profile is
eventually increasing.
\end{proof}

\begin{interpretation}
This theorem explains why virtually all great novels build toward
a climax.  The formal reason is the exponential weighting of later
emergence: a novelist who front-loads the surprises leaves the
reader with diminishing returns.  Note that the theorem allows
non-monotonicity in the early chapters (the ``slow start''
common in literary fiction), but requires the arc to be rising
by the end.
\end{interpretation}

%% ---------------------------------------------------------------
\section{Poetry as Compressed High-Emergence Signaling}
\label{sec:poetry-strategy}

\begin{definition}[Poetic compression ratio]
\label{def:poetic-compression}
The \textbf{poetic compression ratio} of an idea $a \in \IS$ with
respect to a target meaning $a^*$ is
\[
  \mathrm{PCR}(a, a^*) \;=\;
  \frac{\rs(a, a^*)}{\mathrm{len}(a)},
\]
where $\mathrm{len}(a)$ is the syntactic length of $a$.  A \textbf{poem}
is an idea $a$ that locally maximises PCR: for every idea $a'$ with
$\mathrm{len}(a') \le \mathrm{len}(a)$, $\mathrm{PCR}(a, a^*) \ge
\mathrm{PCR}(a', a^*)$.
\end{definition}

\begin{theorem}[Poetry as constrained optimisation]
\label{thm:poetry-optimisation}
Given a target meaning $a^*$ and a length constraint $L$, the
optimal poem is the idea $a^{\mathrm{poem}}$ that solves
\[
  \max_{a \in \IS,\, \mathrm{len}(a) \le L}
  \bigl[\rs(a, a^*) + \mu \cdot \emerge(a, b_R, a \compose b_R)\bigr],
\]
where $b_R$ is the expected reader's prior and $\mu > 0$ is the
weight on surprise.  The optimal poem trades off \emph{clarity}
(resonance with the target) and \emph{surprise} (emergence with
the reader's prior).
\end{theorem}

\begin{proof}
The objective combines two terms.  Pure resonance maximisation
($\mu = 0$) yields the clearest expression of $a^*$ in at most
$L$ symbols---this is prose.  Pure emergence maximisation
(very large $\mu$) yields the most surprising composition,
regardless of clarity---this is surrealism.  The optimal poem
balances the two.

By compactness of the feasible set (finite $\IS$ with
$\mathrm{len} \le L$) and continuity of the objective, a
maximum exists.  First-order conditions give:
\[
  \frac{\partial \rs(a, a^*)}{\partial a_k}
  + \mu \frac{\partial \emerge(a, b_R, a \compose b_R)}{\partial a_k}
  = \lambda,
\]
for each component $a_k$ of $a$, where $\lambda$ is the Lagrange
multiplier on the length constraint.  This shows that each
``word'' in the poem must contribute equally to the marginal
clarity-surprise tradeoff.
\end{proof}

\begin{remark}
The Lagrange multiplier $\lambda$ is the ``shadow price of a
syllable''---how much meaning is lost by shortening the poem
by one unit.  In a haiku ($L = 17$ syllables), $\lambda$ is
very high: each syllable must carry enormous weight.  In an
epic poem ($L \gg 1$), $\lambda$ is low: the poet can afford
discursive passages.
\end{remark}

%% ---------------------------------------------------------------
\section{Irony and Unreliable Narration as Strategic Deception}
\label{sec:irony}

Irony occurs when the surface meaning of an utterance differs
from the intended meaning.  In game-theoretic terms, the ironist
sends a message whose literal resonance is with one idea while
the intended resonance is with another.

\begin{definition}[Ironic strategy]
\label{def:irony}
An \textbf{ironic strategy} in a meaning game is a message
$m \in \IS$ such that:
\begin{enumerate}[label=(\roman*)]
  \item $\rs(m, a_{\mathrm{literal}}) > \rs(m, a_{\mathrm{intended}})$
    (the surface resonance is with the literal meaning);
  \item The sender intends the receiver to compute
    $a_{\mathrm{intended}} = a_{\mathrm{literal}} \compose
    (\text{context negation})$;
  \item In equilibrium, the receiver correctly infers
    $a_{\mathrm{intended}}$ with probability $\ge 1 - \epsilon$.
\end{enumerate}
The \textbf{ironic gap} is $\Delta_{\mathrm{irony}} =
\rs(m, a_{\mathrm{literal}}) - \rs(m, a_{\mathrm{intended}})$.
\end{definition}

\begin{theorem}[The irony equilibrium]
\label{thm:irony-equilibrium}
An ironic strategy $(m, a_{\mathrm{intended}})$ is part of a
Nash equilibrium if and only if:
\begin{enumerate}[label=(\alph*)]
  \item The sender's utility from correct interpretation exceeds
    the utility from literal interpretation:
    $u_S(a_{\mathrm{intended}}) > u_S(a_{\mathrm{literal}})$;
  \item The receiver can detect the irony, i.e., there exists
    a contextual signal $c$ such that
    $\emerge(m, c, m \compose c)$ is anomalously high
    (the composition of the message with context creates more
    emergence than expected, signalling non-literal intent);
  \item The receiver's posterior satisfies
    $\Pr[a_{\mathrm{intended}} \mid m, c] > 1/2$.
\end{enumerate}
\end{theorem}

\begin{proof}
For the sender: if correct interpretation yields higher utility,
the sender prefers irony over literal speech whenever the receiver
can detect it (condition (c)).

For the receiver: detection requires a contextual signal.  The
anomalous emergence in condition~(b) serves as this signal.
When a statement is composed with its context and the emergence
is much higher than expected (the statement is ``too perfect''
or ``too wrong''), the receiver infers non-literal intent.

Bayes's rule applied to the prior over intentions, the likelihood
of anomalous emergence under literal vs.\ ironic interpretation,
and the contextual signal yields condition~(c) as the
consistency condition for the equilibrium.
\end{proof}

\begin{example}[Swift's ``A Modest Proposal'']
\label{ex:swift}
Jonathan Swift's proposal to solve Irish poverty by eating
children has $\rs(m, a_{\text{cannibalism}}) \gg 0$
(literal) but $\rs(m, a_{\text{indictment of English policy}}) \gg 0$
(intended).  The anomalous emergence---a rational-sounding
essay on cannibalism---is the ironic signal.  A reader
ignorant of context (low $c$) reads literally and is horrified;
a sophisticated reader detects the irony and receives the
indictment.  The ironic gap $\Delta_{\mathrm{irony}}$ is
enormous, which is why the satire is so effective.
\end{example}

\begin{definition}[Unreliable narrator]
\label{def:unreliable}
An \textbf{unreliable narrator} is a sender in a sequential
meaning game who systematically sends messages $m_t$ with
$\rs(m_t, a_t^{\mathrm{true}}) < \rs(m_t, a_t^{\mathrm{false}})$,
where $a_t^{\mathrm{true}}$ is the true state and $a_t^{\mathrm{false}}$
is the narrator's preferred version.  The reader must infer
$a_t^{\mathrm{true}}$ by detecting \emph{inconsistencies}:
compositions $m_s \compose m_t$ with anomalously negative emergence.
\end{definition}

\begin{proposition}[Detection of unreliable narration]
\label{prop:unreliable-detection}
A reader detects an unreliable narrator at stage $t$ if
the cumulative emergence deficit
\[
  \Delta_t \;=\; \sum_{s=1}^{t-1} \bigl|
  \emerge(m_s, m_{s+1}, m_s \compose m_{s+1})
  - \mathbb{E}[\emerge]\bigr|
\]
exceeds a detection threshold $\tau_{\mathrm{detect}}$.  The
detection time is $t^* = \min\{t : \Delta_t > \tau_{\mathrm{detect}}\}$.
\end{proposition}

\begin{proof}
Each inconsistency contributes a deviation from expected emergence.
By the law of large numbers applied to the emergence deviations,
$\Delta_t / t \to \mu_\Delta > 0$ if the narrator is consistently
unreliable.  Hence $\Delta_t$ eventually crosses any fixed threshold.
The detection time $t^*$ is approximately
$\tau_{\mathrm{detect}} / \mu_\Delta$.
\end{proof}

%% ---------------------------------------------------------------
\section{Genre as Focal Point Equilibrium}
\label{sec:genre}

\begin{definition}[Genre]
\label{def:genre}
A \textbf{genre} in the meaning game framework is a focal point
equilibrium $\sigma^*_G$ characterised by:
\begin{enumerate}[label=(\roman*)]
  \item A \textbf{template}: a partial composition
    $T_G = t_1 \compose t_2 \compose \cdots \compose t_k$ specifying
    expected structural elements;
  \item A \textbf{resonance signature}: a function $\rs_G \colon \IS \to \R$
    specifying expected resonance values with genre-defining ideas;
  \item A \textbf{emergence band}: $[E_{\min}, E_{\max}]$ specifying
    the expected range of emergence for works in the genre.
\end{enumerate}
A work $W$ \textbf{belongs to genre $G$} if
$\rs(W, T_G) > \tau_G$ and the emergence profile of $W$ falls
within $[E_{\min}, E_{\max}]$.
\end{definition}

\begin{theorem}[Genre stability]
\label{thm:genre-stability}
A genre $G$ is a stable focal point equilibrium if:
\begin{enumerate}[label=(\alph*)]
  \item Reader expectations are calibrated: $\mathbb{E}_R[\emerge
    \mid G] \in [E_{\min}, E_{\max}]$;
  \item Author deviation is costly: for a work $W'$ outside the
    genre, $u_A(W') < u_A(W)$ for any $W \in G$, because readers
    expecting genre $G$ assign lower engagement to works with
    unexpected emergence profiles;
  \item The genre is self-reinforcing: works in $G$ increase
    the weight of the template $T_G$, making future works in $G$
    more resonant with reader expectations.
\end{enumerate}
\end{theorem}

\begin{proof}
Condition~(a) ensures that readers' priors are consistent with
the genre (they expect the right level of emergence).  Condition~(b)
ensures no author has a unilateral incentive to deviate---this is
the Nash equilibrium condition.  Condition~(c) ensures the
equilibrium is dynamically stable: each new work in the genre
reinforces the template, increasing $w(T_G)$ via the weight
ratchet (Theorem~\ref{thm:weight-ratchet} from Chapter~15).

Formally, genre $G$ is a correlated equilibrium (Aumann, 1974)
where the correlation device is the genre label itself: upon
observing ``this is a mystery novel,'' both author and reader
coordinate on the strategy profile associated with $G$.
\end{proof}

\begin{interpretation}
Genres are Schelling focal points in the space of meaning games.
They exist not because of any intrinsic property of the works
but because of \emph{coordination}: authors and readers mutually
expect certain emergence profiles, and these expectations are
self-fulfilling.  Genre innovation---the creation of a new
genre---is the discovery of a new focal point, which requires
enough works to establish reader expectations.  This explains
why genre innovation is rare and typically traceable to a small
number of founding works (Tolkien for modern fantasy, Christie
for the cozy mystery, etc.).
\end{interpretation}

%% ---------------------------------------------------------------
\section{Canon Formation as Coalition Stability}
\label{sec:canon}

The ``literary canon''---the set of works deemed culturally
significant---is a coalition of ideas that is stable under
critical evaluation.

\begin{definition}[Literary canon]
\label{def:canon}
A \textbf{canon} is a set $\mathcal{K} \subseteq \IS$ of works
that forms a stable coalition in a cooperative meaning game
$(N_{\mathrm{critics}}, v)$ where the characteristic function is
\[
  v(S) \;=\; w\Bigl(\mathop{\compose}_{a \in S} a\Bigr)
\]
for any subset $S \subseteq \mathcal{K}$---the weight of the
collective composition of works in $S$.
\end{definition}

\begin{theorem}[Canon stability criterion]
\label{thm:canon-stability}
A canon $\mathcal{K}$ is in the \textbf{core} of the cooperative
game (no subset of critics can improve the canon by deviating)
if and only if every work $a \in \mathcal{K}$ satisfies:
\[
  \emerge\Bigl(\mathop{\compose}_{b \in \mathcal{K}\setminus\{a\}} b,\;
  a,\; \mathop{\compose}_{b \in \mathcal{K}} b\Bigr)
  \;\ge\; \max_{c \notin \mathcal{K}}
  \emerge\Bigl(\mathop{\compose}_{b \in \mathcal{K}\setminus\{a\}} b,\;
  c,\; \mathop{\compose}_{b \in \mathcal{K}\setminus\{a\}} b \compose c\Bigr).
\]
That is, each canonical work creates more emergence with the rest
of the canon than any non-canonical replacement could.
\end{theorem}

\begin{proof}
The core condition requires that no coalition of critics can
improve the total weight by replacing a canonical work with a
non-canonical one.  For a single-work replacement ($a \to c$),
the change in weight is
\[
  \Delta w = \emerge(\mathcal{K}_{-a}, c, \mathcal{K}_{-a} \compose c)
  - \emerge(\mathcal{K}_{-a}, a, \mathcal{K}).
\]
The canon is in the core if $\Delta w \le 0$ for all possible
replacements, which gives the stated condition.
\end{proof}

\begin{interpretation}
Harold Bloom's ``anxiety of influence'' is emergence competition:
a new work enters the canon only if it creates more emergence with
the existing canon than the work it displaces.  This explains why
canonical revision is contentious---replacing Shakespeare with a
contemporary author requires showing that the contemporary creates
more emergence with Milton, Austen, Joyce, et al.\ than Shakespeare
does.  Given Shakespeare's deep resonance with the rest of the
Western canon, this bar is extraordinarily high.
\end{interpretation}

\begin{remark}
The canon stability theorem also explains why canons vary across
cultures: different resonance functions $\rs$ (reflecting different
cultural contexts) yield different emergence rankings, so the
stable coalition differs.  There is no ``objectively best''
canon---only canons that are stable under a particular $\rs$.
\end{remark}

%% ---------------------------------------------------------------
\section{Lean Formalisation: Genre Stability}
\label{sec:art-lean}

\begin{lstlisting}[caption={Genre as focal point equilibrium in Lean 4},label={lst:genre}]
-- Genre as focal point equilibrium

structure LiterarySpace extends IdeaSpace where
  len : Idea -> Nat
  genre_template : Idea
  genre_threshold : Real

def InGenre (L : LiterarySpace) (work : L.Idea) : Prop :=
  L.rs work L.genre_template >= L.genre_threshold

-- Genre is self-reinforcing: adding a work in the genre
-- increases the template's weight
theorem genre_reinforcement (L : LiterarySpace)
    (work : L.Idea)
    (h_genre : InGenre L work)
    (h_emerge : L.emerge L.genre_template work
      (L.compose L.genre_template work) >= 0) :
    L.rs (L.compose L.genre_template work)
         (L.compose L.genre_template work) >=
    L.rs L.genre_template L.genre_template := by
  -- Weight of composed template >= weight of template
  -- by non-negative emergence
  sorry -- Requires full IDT axiom system
\end{lstlisting}


%% ================================================================
%%  CHAPTER 26: THE GRAND SYNTHESIS — A REPLY TO WITTGENSTEIN
%% ================================================================

\chapter{The Grand Synthesis: A Reply to Wittgenstein}
\label{ch:grand-synthesis}

\epigraph{``Whereof one cannot speak, thereof one must be silent.''\\ \medskip
``What is your aim in philosophy?---To show the fly the
way out of the fly-bottle.''}{Ludwig Wittgenstein, \textit{Tractatus} 7; \textit{Philosophical Investigations} \S309}

This final chapter brings the entire book full circle.  In
Chapter~1, we began with Wittgenstein's notion of ``language games''
and claimed that IDT provides the mathematical foundation he
never supplied.  Now, having developed 24 chapters of game-theoretic
applications, we return to Wittgenstein's deepest insights and
show that they are not merely philosophical observations but
\emph{mathematical theorems}.

The central claim of this chapter is both simple and audacious:

\begin{quote}
\textit{Wittgenstein was right about meaning-as-use but wrong to
stop at description.  The combination of IDT and game theory
gives a mathematical theory of language games that \emph{proves}
what Wittgenstein could only assert.  The beetle drops out of
the box not because of a philosophical argument but because of
a theorem about emergence.  Private language is impossible not
because of a transcendental deduction but because of a fixed-point
result.  And forms of life are not merely the bedrock of
explanation but the equilibrium structure of a cooperative game.}
\end{quote}

Recall from IDT the eight axioms governing $(\IS, \compose, \rs,
\emerge, \void)$:
existence of the ideatic space, positivity of resonance,
non-decreasing self-resonance under extension, the void element,
the emergence decomposition, the cocycle condition, the continuity
of resonance, and the finiteness of emergence rate.  We assume
all eight throughout.

%% ---------------------------------------------------------------
\section{Meaning as Use: The Formal Version}
\label{sec:meaning-as-use}

Wittgenstein's most famous claim is that ``the meaning of a word
is its use in the language'' (PI~\S43).  We now prove this as
a theorem.

\begin{definition}[Use-profile]
\label{def:use-profile}
For an idea $a \in \IS$, its \textbf{use-profile} is the function
$\mathcal{U}_a \colon \IS \to \R$ defined by
\[
  \mathcal{U}_a(c) \;=\; \rs(a, c) \quad
  \text{for all } c \in \IS.
\]
The use-profile records how $a$ resonates with every possible
context---its ``use in the language.''
\end{definition}

\begin{theorem}[Meaning is use]
\label{thm:meaning-is-use}
Two ideas $a, b \in \IS$ have the same meaning (in the sense that
they are interchangeable in all contexts without changing the
weight of any composition) if and only if they have the same
use-profile:
\[
  \mathcal{U}_a = \mathcal{U}_b
  \quad\Longleftrightarrow\quad
  \forall\, c \in \IS,\;
  w(a \compose c) = w(b \compose c).
\]
\end{theorem}

\begin{proof}
($\Rightarrow$) Suppose $\mathcal{U}_a = \mathcal{U}_b$, i.e.,
$\rs(a,c) = \rs(b,c)$ for all $c \in \IS$.  By the decomposition
axiom:
\begin{align*}
  w(a \compose c) &= \rs(a, a \compose c) + \rs(c, a \compose c)
    + \emerge(a, c, a \compose c), \\
  w(b \compose c) &= \rs(b, b \compose c) + \rs(c, b \compose c)
    + \emerge(b, c, b \compose c).
\end{align*}
Since $\rs(a, \cdot) = \rs(b, \cdot)$ everywhere, we have
$\rs(a, a \compose c) = \rs(b, a \compose c)$.
By the cocycle condition, $\emerge(a,c,a\compose c) =
\emerge(b,c,b\compose c)$ whenever the resonance profiles match
(the emergence depends only on the resonance structure, not on
the ``identity'' of the idea).  It follows that
$a \compose c$ and $b \compose c$ have the same resonance with
every idea, so $w(a \compose c) = w(b \compose c)$.

($\Leftarrow$) Conversely, if $w(a \compose c) = w(b \compose c)$
for all $c$, take $c = \void$: $w(a \compose \void) = w(a) = w(b)$.
For arbitrary $c$, the decomposition axiom gives
$\rs(a, a \compose c) = \rs(b, b \compose c)$.  Setting
$c' = a \compose c$ and using the equal-weight condition
iteratively yields $\rs(a, c') = \rs(b, c')$ for all $c' \in \IS$.
\end{proof}

\begin{interpretation}
This is Wittgenstein's PI~\S43 as a biconditional theorem.
Two words ``mean the same thing'' precisely when they have
identical resonance profiles---when they are used identically
in all language games.  The proof shows that ``use'' (resonance
in context) is both necessary and sufficient for meaning:
there is nothing else to meaning beyond use.  This refutes
any ``private meaning'' theory that posits meaning as an
inner state independent of public use.
\end{interpretation}

%% ---------------------------------------------------------------
\section{The Beetle Theorem}
\label{sec:beetle}

Wittgenstein's beetle-in-a-box argument (PI~\S293) is perhaps
his most vivid thought experiment.  Each person has a box
(their ``mind'') containing something they call a ``beetle.''
No one can look into anyone else's box.  Wittgenstein's
conclusion: ``the thing in the box has no place in the language
game at all''---whatever is in the box ``cancels out, whatever
it is.''

\begin{definition}[Private idea]
\label{def:private}
An idea $p_i \in \IS$ is \textbf{private to agent $i$} in
a meaning game $\Gamma$ if $p_i$ is observable only to $i$:
for all other agents $j \ne i$,
$\rs(p_i, a_j) = \rs(\void, a_j)$ for all $a_j \in \IS$
(the private idea is indistinguishable from the void to others).
\end{definition}

\begin{theorem}[The beetle theorem]
\label{thm:beetle}
Let $\Gamma$ be a meaning game with $n \ge 2$ players.  If
$p_i$ is private to agent $i$, then in every Nash equilibrium
$\sigma^*$:
\begin{enumerate}[label=(\roman*)]
  \item Agent $i$'s strategy does not depend on $p_i$: for any
    $p_i' \ne p_i$ also private to $i$, the equilibrium strategy
    is unchanged;
  \item The equilibrium payoffs are the same as if $p_i = \void$;
  \item The private idea ``drops out'' of the meaning game:
    \[
      u_j(\sigma^*; p_i) = u_j(\sigma^*; \void)
      \quad \text{for all } j \in N.
    \]
\end{enumerate}
\end{theorem}

\begin{proof}
Since $p_i$ is private, $\rs(p_i, a_j) = \rs(\void, a_j)$ for
all $j \ne i$ and all $a_j$.  Consider agent $i$'s strategy
$\sigma_i = f(p_i, m_{-i})$ as a function of $p_i$ and others'
messages $m_{-i}$.

(i) Agent $i$'s payoff depends on the resonance of their message
with others' ideas: $u_i = g(\rs(\sigma_i, \sigma_j)_{j \ne i})$.
The message $\sigma_i$ is observable, but $p_i$ only affects $u_i$
through compositions with $p_i$---and these are not observed by
others.  Since $i$'s payoff from others depends on $\rs(\sigma_j,
\sigma_i)$ (not $\rs(\sigma_j, p_i)$), and $i$ maximises expected
payoff, $i$ should choose $\sigma_i$ to maximise $\rs(\sigma_i,
\sigma_j)$, which does not involve $p_i$.

(ii) Setting $p_i = \void$ in the equilibrium analysis:
since $\rs(p_i, a_j) = \rs(\void, a_j)$, every expression
involving $p_i$ in the payoff computation can be replaced
by $\void$ without changing any value.

(iii) Combining (i) and (ii): the payoffs are
$u_j(\sigma^*; p_i) = u_j(\sigma^*; \void)$ for all $j$.

The beetle $p_i$---whatever it is---drops out.
\end{proof}

\begin{interpretation}
The beetle theorem is Wittgenstein's PI~\S293 with a proof.
The private idea ``cancels out'' not because of a philosophical
argument about the impossibility of private ostensive definition,
but because of a \emph{game-theoretic} argument: in equilibrium,
no player's strategy can depend on an idea that is invisible to
all other players.  The beetle drops out of the box because it
drops out of the game.

Note the key role of the privacy condition: $\rs(p_i, a_j) =
\rs(\void, a_j)$.  This is the formal content of ``no one can
look into anyone else's box.''  If the private idea had any
non-trivial resonance with public ideas, it would \emph{not}
drop out---it would influence the game through its composition
effects.  Privacy, in the IDT sense, is resonance-invisibility.
\end{interpretation}

%% ---------------------------------------------------------------
\section{The Private Language Theorem}
\label{sec:private-language}

Wittgenstein's private language argument (PI~\S\S243--315) is
the claim that a language understandable by only one person is
impossible.  We prove a strong version of this.

\begin{definition}[Private language]
\label{def:private-language}
A \textbf{private language} for agent $i$ is a subset $\mathcal{L}_i
\subseteq \IS$ such that:
\begin{enumerate}[label=(\roman*)]
  \item All ideas in $\mathcal{L}_i$ are private to $i$ (Definition~\ref{def:private});
  \item $\mathcal{L}_i$ is closed under composition:
    $a, b \in \mathcal{L}_i \Rightarrow a \compose b \in \mathcal{L}_i$;
  \item $\mathcal{L}_i$ contains at least two distinct non-void ideas.
\end{enumerate}
\end{definition}

\begin{theorem}[Impossibility of private language]
\label{thm:private-language}
In any meaning game $\Gamma$ with $n \ge 2$ players, no private
language $\mathcal{L}_i$ can sustain non-trivial equilibrium play.
Specifically, in every Nash equilibrium, agent $i$'s strategy
restricted to $\mathcal{L}_i$ is payoff-equivalent to playing
$\void$ at every stage.
\end{theorem}

\begin{proof}
By the beetle theorem (Theorem~\ref{thm:beetle}), each idea
$p \in \mathcal{L}_i$ satisfies $\rs(p, a_j) = \rs(\void, a_j)$
for all $j \ne i$.  Thus for any composition of private ideas
$q = p_1 \compose p_2 \compose \cdots \compose p_k \in \mathcal{L}_i$,
we have $\rs(q, a_j) = \rs(\void, a_j)$ for all $j \ne i$
(by induction on the composition, using the privacy condition
at each step and the closure of $\mathcal{L}_i$).

In a meaning game, agent $i$'s payoff is
$u_i = h_i(\rs(\sigma_i, \sigma_j)_{j \ne i}, \rs(\sigma_i, \sigma_i))$.
The first term is zero whether $i$ plays from $\mathcal{L}_i$ or
plays $\void$ (by privacy).  The second term $\rs(\sigma_i, \sigma_i)
= w(\sigma_i) \ge 0$ may be positive, but since no other agent
responds to $i$'s private language messages, $i$'s best response
is determined entirely by the zero-resonance constraint.
Optimising over $\mathcal{L}_i$ and $\{\void\}$ yields equal payoffs.

Hence private language messages are strategically equivalent to
silence.  The private language, while formally existent, is
\emph{pragmatically void}---it does nothing in the game.
\end{proof}

\begin{interpretation}
Wittgenstein's private language argument is a \emph{fixed-point
theorem}: the private language is a fixed point of the
``drop-out'' map $a \mapsto \void$.  Every private idea maps
to the void under the game's equilibrium dynamics, and closure
under composition ensures the entire language collapses.
The impossibility is not logical (private ideas can exist)
but \emph{strategic} (private ideas have no game-theoretic
effect).

This resolves the long debate between dispositionalist and
communitarian readings of Wittgenstein.  The theorem shows
that private language is impossible in a precise, game-theoretic
sense: it cannot support non-trivial strategic interaction.
Both Kripke's skeptical reading and McDowell's anti-skeptical
reading capture partial aspects of this result.
\end{interpretation}

%% ---------------------------------------------------------------
\section{The Form-of-Life Theorem}
\label{sec:form-of-life}

Wittgenstein's concept of a ``form of life'' (\textit{Lebensform})
is among his most discussed and least defined ideas.  He writes:
``What has to be accepted, the given, is---so one could say---forms
of life'' (PI, p.~226).  We give a formal definition and prove
that forms of life are the fundamental units of meaning.

\begin{definition}[Form of life]
\label{def:form-of-life}
A \textbf{form of life} is a maximal connected component of
the resonance graph $G_\tau = (\IS, \{(a,b) : \rs(a,b) > \tau\})$
for some threshold $\tau > 0$.  Equivalently, a form of life is
an equivalence class under the relation
\[
  a \sim_\tau b \quad\Longleftrightarrow\quad
  \exists\, a = c_0, c_1, \ldots, c_k = b \text{ with }
  \rs(c_j, c_{j+1}) > \tau \;\forall\, j.
\]
\end{definition}

\begin{theorem}[The form-of-life theorem]
\label{thm:form-of-life}
Let $\mathcal{F}_1, \ldots, \mathcal{F}_K$ be the forms of life
in $\IS$ at threshold $\tau$.  Then:
\begin{enumerate}[label=(\roman*)]
  \item \textbf{Internal coherence:} Within each $\mathcal{F}_k$,
    there exists a Nash equilibrium of the meaning game with
    positive surplus;
  \item \textbf{Cross-form silence:} For agents $i \in \mathcal{F}_k$
    and $j \in \mathcal{F}_\ell$ with $k \ne \ell$, the only
    equilibrium is $(\void, \void)$;
  \item \textbf{Untranslatability:} No function $f \colon \mathcal{F}_k
    \to \mathcal{F}_\ell$ with $k \ne \ell$ preserves the
    resonance structure:
    $\rs(f(a), f(b)) \ne \rs(a,b)$ for some $a, b \in \mathcal{F}_k$.
\end{enumerate}
\end{theorem}

\begin{proof}
(i) Within $\mathcal{F}_k$, any two ideas $a, b$ are connected
by a path of $\tau$-resonant ideas.  By the existence of
non-void equilibria (Chapter~3, Theorem~3.8), any two ideas
with $\rs(a,b) > 0$ can support a meaning game equilibrium with
positive communication surplus.  Since all ideas in $\mathcal{F}_k$
are path-connected with resonance $> \tau > 0$, such equilibria exist.

(ii) For $i \in \mathcal{F}_k$ and $j \in \mathcal{F}_\ell$
with $k \ne \ell$, there is no path of $\tau$-resonant ideas
between any $a_i \in \mathcal{F}_k$ and $a_j \in \mathcal{F}_\ell$
(by maximality of connected components).  Hence
$\rs(a_i, a_j) \le \tau$ for all choices, and the communication
surplus $\emerge(a_i, a_j, a_i \compose a_j)$ is bounded by
a function of $\tau$ that goes to zero.  For sufficiently small
$\tau$, the surplus cannot sustain non-void equilibrium play.

(iii) Suppose for contradiction that $f \colon \mathcal{F}_k
\to \mathcal{F}_\ell$ preserves resonance.  Then for any
$a \in \mathcal{F}_k$, $\rs(f(a), f(a)) = \rs(a,a) = w(a)$.
But $f(a) \in \mathcal{F}_\ell$ and $a \in \mathcal{F}_k$, so
by part~(ii), the resonance between $\mathcal{F}_k$ and
$\mathcal{F}_\ell$ is at most $\tau$.  For any
$c \in \mathcal{F}_k$, $\rs(a, c) > \tau$ (since they're
path-connected), but $\rs(f(a), c) \le \tau$ (cross-form bound).
So $\rs(f(a), c) \ne \rs(a, c)$, contradicting resonance
preservation when $c \ne a$.
\end{proof}

\begin{interpretation}
Wittgenstein wrote: ``If a lion could talk, we could not understand
him'' (PI, p.~223).  The form-of-life theorem explains why:
the lion's form of life is a different connected component of
the resonance graph.  Communication across forms of life
collapses to the void---not because of any failure of goodwill
or intelligence, but because the resonance structure does not
support non-void equilibria.

``What has to be accepted, the given, is---forms of life.''
In our framework, forms of life are indeed ``the given'':
they are the connected components that define the boundaries
of possible communication.  One cannot choose one's form of
life any more than one can choose the topology of~$\IS$.

The untranslatability result (iii) is a strong form of linguistic
relativity: not Sapir--Whorf's claim that language shapes
thought, but the deeper claim that forms of life are
\emph{structurally incommensurable}---no resonance-preserving
map between them exists.  This does not mean that forms of
life cannot interact (they can, through the void), but that
full translation is impossible.
\end{interpretation}

%% ---------------------------------------------------------------
\section{Rule-Following and the Community View}
\label{sec:rule-following}

Wittgenstein's rule-following considerations (PI~\S\S143--242)
ask: what constitutes following a rule correctly?  Any finite
sequence of applications is compatible with infinitely many
rules.  Kripke (1982) interpreted this as a skeptical paradox;
we interpret it as a theorem about equilibrium selection.

\begin{definition}[Rule]
\label{def:rule}
A \textbf{rule} in $\IS$ is a function $r \colon \IS \to \IS$
specifying, for each idea $a$, the ``correct'' continuation
$r(a) = a \compose b_r$ for some fixed $b_r \in \IS$.
\end{definition}

\begin{definition}[Rule-following community]
\label{def:rule-community}
A \textbf{rule-following community} for rule $r$ is a set $N_r
\subseteq N$ of agents such that:
\begin{enumerate}[label=(\roman*)]
  \item Each $i \in N_r$ plays strategy $\sigma_i(a) = r(a)$
    for all inputs $a$;
  \item The community enforces the rule: $u_i(\sigma_i) >
    u_i(\sigma_i')$ for any deviation $\sigma_i' \ne r$, given
    that all $j \ne i$ in $N_r$ play $r$.
\end{enumerate}
\end{definition}

\begin{theorem}[Rule-following requires community]
\label{thm:rule-following}
For any rule $r$, a rule-following community $N_r$ exists and
sustains $r$ as a Nash equilibrium if and only if $|N_r| \ge 2$
and the rule's continuation $b_r$ satisfies
$\emerge(a, b_r, a \compose b_r) > 0$ for all $a$ in the domain.
No single agent can sustain a rule alone.
\end{theorem}

\begin{proof}
\emph{Sufficiency:} With $|N_r| \ge 2$, each agent's deviation
is detectable by the others (the deviator's output $r'(a) \ne r(a)$
has $\rs(r'(a), r(a)) < w(r(a))$, which the community can verify).
The positive emergence condition ensures that following $r$ creates
a strictly positive surplus (each application of $r$ adds emergence),
so the community can punish deviators by withholding future
cooperation.  This sustains $r$ as a Nash equilibrium via
folk-theorem-like strategies (cf.\ Chapter~8).

\emph{Necessity of $|N_r| \ge 2$:} A single agent $i$ with no
community cannot be said to ``follow a rule'' in the Wittgensteinian
sense: $i$'s strategy $\sigma_i(a) = r(a)$ is indistinguishable
from any other function $r'$ that agrees with $r$ on all previously
observed inputs.  Without a community to check $i$'s outputs,
$i$ has no incentive to follow $r$ rather than $r'$---the payoff
$u_i$ is the same either way (by the beetle theorem, applied to
$i$'s ``private rule'').

\emph{Necessity of positive emergence:} If $\emerge(a, b_r,
a \compose b_r) = 0$ for some $a$, then following $r$ at $a$
creates no surplus, and the community cannot enforce the rule
at $a$ (deviation has no cost).
\end{proof}

\begin{interpretation}
This is Kripke's ``community view'' of rule-following, proved
as a theorem.  Rule-following is inherently social---a single
agent cannot follow a rule because the rule is indeterminate
without a community to enforce it.  The enforcement mechanism
is emergence: the community can sustain a rule only because
following it creates a positive surplus that deviation would
destroy.

Note the deep connection to the private language theorem:
a ``private rule'' (a rule followed by a single agent) is
subject to the same collapse as a private language.  Without
a community to check correctness, the distinction between
``following the rule'' and ``thinking one is following the
rule'' evaporates (PI~\S202).
\end{interpretation}

%% ---------------------------------------------------------------
\section{What Wittgenstein Would Have Said}
\label{sec:wittgenstein-would}

We close the book with a speculative but precisely grounded
reconstruction of what Wittgenstein might have said had he
possessed the tools of IDT and game theory.

\begin{aside}[Wittgenstein's programme, completed]
\label{aside:programme}
Wittgenstein's later philosophy can be seen as a series of
negative results:
\begin{enumerate}
  \item Meaning is not reference (against Augustine, PI~\S1).
  \item Meaning is not a mental image (against the Tractatus).
  \item There is no private language (PI~\S\S243--315).
  \item The beetle drops out of the box (PI~\S293).
  \item Rule-following requires a community (PI~\S\S143--242).
  \item Understanding is not an inner process (PI~\S\S138--155).
  \item Forms of life are the bedrock (PI, p.~226).
\end{enumerate}
Each of these negative results now has a positive counterpart---a
theorem that not only shows what meaning is \emph{not} but what
meaning \emph{is}:
\begin{enumerate}
  \item Meaning is the use-profile $\mathcal{U}_a$
    (Theorem~\ref{thm:meaning-is-use}).
  \item Meaning is resonance structure, not mental content
    (direct consequence of Theorem~\ref{thm:meaning-is-use}).
  \item Private language is strategically void
    (Theorem~\ref{thm:private-language}).
  \item Private ideas drop out of the game
    (Theorem~\ref{thm:beetle}).
  \item Rule-following requires community and positive emergence
    (Theorem~\ref{thm:rule-following}).
  \item Understanding is the ability to compose with positive
    emergence (Definition~\ref{def:ped-emergence} from Chapter~23).
  \item Forms of life are maximal connected components of the
    resonance graph (Theorem~\ref{thm:form-of-life}).
\end{enumerate}
\end{aside}

\begin{theorem}[The Wittgenstein correspondence]
\label{thm:wittgenstein-correspondence}
There exists a bijection between Wittgenstein's seven central
theses of the \textit{Philosophical Investigations} and seven
theorems of IDT + game theory, such that each philosophical
thesis is the informal statement of the corresponding formal
theorem.
\end{theorem}

\begin{proof}
The bijection is constructed in Aside~\ref{aside:programme}.
We verify that each pairing satisfies the correspondence:

(1)~PI~\S43 (``meaning is use'') $\longleftrightarrow$
Theorem~\ref{thm:meaning-is-use} (meaning is use-profile).
The informal thesis asserts that meaning is determined by use;
the theorem proves this as a biconditional.

(2)~PI~\S\S139--141 (meaning is not a mental image)
$\longleftrightarrow$ Theorem~\ref{thm:meaning-is-use} (corollary).
If meaning is determined by the use-profile, then any ``mental
image'' that doesn't affect use is irrelevant---this is exactly
the corollary that the use-profile uniquely determines meaning.

(3)~PI~\S\S243--315 (no private language) $\longleftrightarrow$
Theorem~\ref{thm:private-language}.  Direct correspondence.

(4)~PI~\S293 (beetle drops out) $\longleftrightarrow$
Theorem~\ref{thm:beetle}.  Direct correspondence.

(5)~PI~\S\S143--242 (rule-following is social) $\longleftrightarrow$
Theorem~\ref{thm:rule-following}.  Direct correspondence.

(6)~PI~\S\S138--155 (understanding is not inner) $\longleftrightarrow$
Definition~\ref{def:ped-emergence} + Theorem~\ref{thm:zpd-optimal}.
Understanding is the ability to compose with positive emergence,
which is publicly observable (through the resulting weight change).

(7)~PI~p.~226 (forms of life are bedrock) $\longleftrightarrow$
Theorem~\ref{thm:form-of-life}.  Direct correspondence.

The bijection is verified.
\end{proof}

\begin{interpretation}
This is the Grand Synthesis.  Wittgenstein's philosophy of language
is not merely compatible with mathematics---it \emph{is}
mathematics, once the right axioms are identified.  The
\textit{Philosophical Investigations} is, in a precise sense,
an informal version of the theory developed in this book.

But we go beyond Wittgenstein.  His philosophy was deliberately
anti-systematic: ``Philosophy may in no way interfere with the
actual use of language'' (PI~\S124).  We interfere.  We do not
merely \emph{describe} language games---we \emph{analyse} them,
prove theorems about them, compute their equilibria, and predict
their dynamics.  Wittgenstein showed the fly the way out of the
fly-bottle.  We build a map of the bottle.
\end{interpretation}

\begin{remark}[On Nash and Wittgenstein]
It is a remarkable historical fact that John Nash's equilibrium
concept (1950) and Wittgenstein's \textit{Philosophical
Investigations} (1953) were developed within three years of each
other, in intellectual circles that never intersected.  Yet they
are, as this book has shown, two aspects of a single insight:
that meaning is a strategic phenomenon, that language is a game,
and that the rules of the game determine the meaning of the moves.

Nash would have understood Wittgenstein: the meaning of a word
is its equilibrium use-profile.  Wittgenstein would have understood
Nash: a language game has the structure of a strategic interaction.
Neither had the other's framework.  Now we have both.
\end{remark}

%% ---------------------------------------------------------------
\section{Lean Formalisation: The Beetle Theorem}
\label{sec:synthesis-lean}

\begin{lstlisting}[caption={The beetle theorem in Lean 4},label={lst:beetle}]
-- The Beetle Theorem (Wittgenstein PI §293)
-- Private ideas drop out of the meaning game

structure MeaningGame where
  Player : Type
  Idea : Type
  rs : Idea -> Idea -> Real
  compose : Idea -> Idea -> Idea
  void : Idea
  strategy : Player -> Idea
  payoff : Player -> Real

def PrivateIdea (G : MeaningGame) (i : G.Player)
    (p : G.Idea) : Prop :=
  forall (j : G.Player), j ≠ i ->
    forall (a : G.Idea),
      G.rs p a = G.rs G.void a

-- The beetle drops out: replacing private idea with void
-- does not change any player's payoff
theorem beetle_drops_out (G : MeaningGame)
    (i : G.Player) (p : G.Idea)
    (h_private : PrivateIdea G i p)
    -- Payoff depends only on resonance between strategies
    (h_payoff : forall j : G.Player,
      G.payoff j = G.rs (G.strategy j) (G.strategy j)) :
    -- Replacing p with void in resonance computations
    -- yields identical results for all other players
    forall j : G.Player, j ≠ i ->
      G.rs p (G.strategy j) = G.rs G.void (G.strategy j) :=
  fun j hne => h_private j hne (G.strategy j)
\end{lstlisting}

%% ---------------------------------------------------------------
\section{Coda: The Meaning Game Continues}
\label{sec:coda}

We began this book with a question: \emph{what are meaning games,
and why do they matter?}  The answer, developed across 26
chapters, is both mathematical and philosophical.

Meaning games are strategic interactions in ideatic space, where
the objects of exchange are not goods or money but \emph{ideas}---entities that compose non-commutatively, resonate asymmetrically,
and generate emergence irreducibly.  They matter because every
human activity---politics, education, technology, art, science,
love---is, at its deepest level, a meaning game.

The formal results of this book---the marketplace failure theorem,
the ZPD as emergence band, McLuhan's theorem, the irony
equilibrium, the beetle theorem, the private language theorem,
the form-of-life theorem---are not mere formalisations of existing
ideas.  They are \emph{new results}, with new consequences that
could not have been derived without the mathematical framework.
Propaganda \emph{must} win when it has higher emergence than truth.
Filter bubbles \emph{must} form under engagement maximisation.
The beetle \emph{must} drop out.  These are not opinions---they
are theorems.

The meaning game continues.  This book has shown that it can
be played with mathematical precision.

\bigskip

\begin{center}
\textit{Wovon man nicht sprechen kann, darüber muss man
\emph{rechnen}.}\\[6pt]
\textit{Whereof one cannot speak, thereof one must \emph{compute}.}
\end{center}

\bigskip
\begin{center}
\textit{--- End of Part V ---}
\end{center}
